% GENERATED FILE. Edit canonical lessons, then run npm run generate:books.
% canonical-source-hash: fnv1a64:cefe551e3d9614a6
% canonical-lessons: MR-C389-amsatah, MR-C389-purnapane, MR-C389-keval, MR-C389-nivval, MR-C389-sahsa

\chapter{How Much and How Often}
\label{ch:mr-how-much-and-how-often}
\hlchaptermodality{389}

\begin{chapteropening}
\textbf{By the end of this chapter:} \emph{I can say partly, completely, merely and usually.}

Complete the last lesson of chapter 389: I can say partly, completely, merely and usually.
\end{chapteropening}

\section[\texorpdfstring{\mr{अंशतः} (aṁśataḥ)}{अंशतः (aṁśataḥ)}]{\mr{अंशतः} (aṁśataḥ) — partly}

\label{lesson:MR-C389-amsatah}



\begin{quote}
Before the new one: say the Marathi for development, then the Marathi for to implement.
\end{quote}

\subsection*{You'll want to know: \mr{अंशतः}}
\textbf{\mr{अंशतः}} — \emph{aṁśataḥ} — \textquotedblleft{}partly\textquotedblright{}.

\begin{quote}
\mr{तुमचं} \mr{म्हणणं} \mr{अंशतः} \mr{खरं} \mr{आहे}.
\end{quote}

\begin{quote}
\emph{tumcaṁ mhaṇṇaṁ aṁśataḥ kharaṁ āhe.}
\end{quote}

\begin{quote}
What you say is partly true.
\end{quote}

Formal; \mr{अंश}, a part.

\begin{cousinweb}
\textbf{\mr{अंशतः}} (\emph{aṁśataḥ}) is a learned adverb built from Sanskrit \emph{aṁśa}, \textquotedblleft{}share, part\textquotedblright{} (CDIAL 2), with the ending \emph{-taḥ}, \textquotedblleft{}by, in\textquotedblright{}.
\end{cousinweb}

\begin{grammarlens}[title={What you've built}]
One more word for saying what happens if.
\end{grammarlens}

\subsection*{Your turn}
\begin{itemize}
\raggedright
  \item \emph{Say it:} \emph{aṁśataḥ}
  \item \emph{Say it:} \emph{aṁśataḥ}, once more
  \item \emph{Say it:} \emph{tumcaṁ mhaṇṇaṁ aṁśataḥ kharaṁ āhe.} — What you say is partly true.
  \item \emph{Recall:} say the Marathi for development, then the Marathi for to implement, then say \emph{aṁśataḥ} again
\end{itemize}

\begin{tcolorbox}[breakable,colback=teal!4,colframe=teal!35!black,title={Before you move on}]
What does \mr{अंशतः} mean? (\textquotedblleft{}Partly\textquotedblright{}.) Say it once more.
\end{tcolorbox}
\section[\texorpdfstring{\mr{पूर्णपणे} (pūrṇapaṇe)}{पूर्णपणे (pūrṇapaṇe)}]{\mr{पूर्णपणे} (pūrṇapaṇe) — completely}

\label{lesson:MR-C389-purnapane}



\begin{quote}
Before the new one: say the Marathi for to implement, then the Marathi for partly.
\end{quote}

\subsection*{You'll want to know: \mr{पूर्णपणे}}
\textbf{\mr{पूर्णपणे}} — \emph{pūrṇapaṇe} — \textquotedblleft{}completely\textquotedblright{}.

\begin{quote}
\mr{पावसामुळे} \mr{रस्ता} \mr{पूर्णपणे} \mr{बंद} \mr{झाला}.
\end{quote}

\begin{quote}
\emph{pāvsāmuḷe rastā pūrṇapaṇe band jhālā.}
\end{quote}

\begin{quote}
Because of the rain the road was completely closed.
\end{quote}

\mr{पूर्ण}, full, + -\mr{पणे}, the adverb ending.

\begin{cousinweb}
\textbf{\mr{पूर्णपणे}} (\emph{pūrṇapaṇe}) is built from \mr{पूर्ण}, \textquotedblleft{}full\textquotedblright{}, a tatsama from Sanskrit \emph{pūrṇa}, \textquotedblleft{}filled, full\textquotedblright{} (CDIAL 8339), with the adverb ending -\mr{पणे}.
\end{cousinweb}

\begin{grammarlens}[title={What you've built}]
One more word for saying what happens if.
\end{grammarlens}

\subsection*{Your turn}
\begin{itemize}
\raggedright
  \item \emph{Say it:} \emph{pūrṇapaṇe}
  \item \emph{Say it:} \emph{pūrṇapaṇe}, once more
  \item \emph{Say it:} \emph{pāvsāmuḷe rastā pūrṇapaṇe band jhālā.} — Because of the rain the road was completely closed.
  \item \emph{Recall:} say the Marathi for to implement, then the Marathi for partly, then say \emph{pūrṇapaṇe} again
\end{itemize}

\begin{tcolorbox}[breakable,colback=teal!4,colframe=teal!35!black,title={Before you move on}]
What does \mr{पूर्णपणे} mean? (\textquotedblleft{}Completely\textquotedblright{}.) Say it once more.
\end{tcolorbox}
\section[\texorpdfstring{\mr{केवळ} (kevaḷ)}{केवळ (kevaḷ)}]{\mr{केवळ} (kevaḷ) — merely; only}

\label{lesson:MR-C389-keval}



\begin{quote}
Before the new one: say the Marathi for partly, then the Marathi for completely.
\end{quote}

\subsection*{You'll want to know: \mr{केवळ}}
\textbf{\mr{केवळ}} — \emph{kevaḷ} — \textquotedblleft{}merely; only\textquotedblright{}.

\begin{quote}
\mr{हे} \mr{काम} \mr{केवळ} \mr{एका} \mr{दिवसात} \mr{झालं}.
\end{quote}

\begin{quote}
\emph{he kām kevaḷ ekā divsāt jhālaṁ.}
\end{quote}

\begin{quote}
This work got done in only one day.
\end{quote}

Formal for \mr{फक्त}.

\begin{cousinweb}
\textbf{\mr{केवळ}} (\emph{kevaḷ}) is a tatsama, taken from Sanskrit \emph{kevala}, \textquotedblleft{}alone, sole, entire\textquotedblright{} (Monier-Williams).
\end{cousinweb}

\begin{grammarlens}[title={What you've built}]
One more word for saying what happens if.
\end{grammarlens}

\subsection*{Your turn}
\begin{itemize}
\raggedright
  \item \emph{Say it:} \emph{kevaḷ}
  \item \emph{Say it:} \emph{kevaḷ}, once more
  \item \emph{Say it:} \emph{he kām kevaḷ ekā divsāt jhālaṁ.} — This work got done in only one day.
  \item \emph{Recall:} say the Marathi for partly, then the Marathi for completely, then say \emph{kevaḷ} again
\end{itemize}

\begin{tcolorbox}[breakable,colback=teal!4,colframe=teal!35!black,title={Before you move on}]
What does \mr{केवळ} mean? (\textquotedblleft{}Merely; only\textquotedblright{}.) Say it once more.
\end{tcolorbox}
\section[\texorpdfstring{\mr{निव्वळ} (nivvaḷ)}{निव्वळ (nivvaḷ)}]{\mr{निव्वळ} (nivvaḷ) — purely; sheer}

\label{lesson:MR-C389-nivval}



\begin{quote}
Before the new one: say the Marathi for completely, then the Marathi for merely.
\end{quote}

\subsection*{You'll want to know: \mr{निव्वळ}}
\textbf{\mr{निव्वळ}} — \emph{nivvaḷ} — \textquotedblleft{}purely; sheer\textquotedblright{}.

\begin{quote}
\mr{हा} \mr{निव्वळ} \mr{वेडेपणा} \mr{आहे}.
\end{quote}

\begin{quote}
\emph{hā nivvaḷ veḍepṇā āhe.}
\end{quote}

\begin{quote}
This is sheer madness.
\end{quote}

Before a noun, sheer: \mr{निव्वळ} \mr{खोटं}, an outright lie.

\begin{cousinweb}
\textbf{\mr{निव्वळ}} (\emph{nivvaḷ}) is a Marathi word whose origin is not settled.
\end{cousinweb}

\begin{grammarlens}[title={What you've built}]
One more word for saying what happens if.
\end{grammarlens}

\subsection*{Your turn}
\begin{itemize}
\raggedright
  \item \emph{Say it:} \emph{nivvaḷ}
  \item \emph{Say it:} \emph{nivvaḷ}, once more
  \item \emph{Say it:} \emph{hā nivvaḷ veḍepṇā āhe.} — This is sheer madness.
  \item \emph{Recall:} say the Marathi for completely, then the Marathi for merely, then say \emph{nivvaḷ} again
\end{itemize}

\begin{tcolorbox}[breakable,colback=teal!4,colframe=teal!35!black,title={Before you move on}]
What does \mr{निव्वळ} mean? (\textquotedblleft{}Purely; sheer\textquotedblright{}.) Say it once more.
\end{tcolorbox}
\section[\texorpdfstring{\mr{सहसा} (sahsā)}{सहसा (sahsā)}]{\mr{सहसा} (sahsā) — usually}

\label{lesson:MR-C389-sahsa}



\begin{quote}
Before the new one: say the Marathi for merely, then the Marathi for purely.
\end{quote}

\subsection*{You'll want to know: \mr{सहसा}}
\textbf{\mr{सहसा}} — \emph{sahsā} — \textquotedblleft{}usually\textquotedblright{}.

\begin{quote}
\mr{मी} \mr{सहसा} \mr{सकाळी} \mr{लवकर} \mr{उठतो}.
\end{quote}

\begin{quote}
\emph{mī sahsā sakāḷī lavkar uṭhto.}
\end{quote}

\begin{quote}
I usually get up early in the morning.
\end{quote}

With a negative, not readily: \mr{तो} \mr{सहसा} \mr{रागावत} \mr{नाही}.

\begin{cousinweb}
\textbf{\mr{सहसा}} (\emph{sahsā}) is usually traced to Sanskrit \emph{sahasā}, \textquotedblleft{}forcibly, suddenly\textquotedblright{} (Monier-Williams); in Marathi it means as a rule.
\end{cousinweb}

\begin{grammarlens}[title={What you've built}]
One more word for saying what happens if.
\end{grammarlens}

\subsection*{Your turn}
\begin{itemize}
\raggedright
  \item \emph{Say it:} \emph{sahsā}
  \item \emph{Say it:} \emph{sahsā}, once more
  \item \emph{Say it:} \emph{mī sahsā sakāḷī lavkar uṭhto.} — I usually get up early in the morning.
  \item \emph{Recall:} say the Marathi for merely, then the Marathi for purely, then say \emph{sahsā} again
\end{itemize}

\begin{tcolorbox}[breakable,colback=teal!4,colframe=teal!35!black,title={Before you move on}]
What does \mr{सहसा} mean? (\textquotedblleft{}Usually\textquotedblright{}.) Say it once more.
\end{tcolorbox}
